% Figure: Diophantine solutions as equally spaced lattice points on a line.
% Source for Figure~\ref{fig:divline}. TikZ only; regenerable by compiling main.tex.
\begin{tikzpicture}[
  pt/.style={circle,draw,inner sep=1.6pt,minimum size=3mm},
  sol/.style={circle,draw,fill=accent!30,inner sep=1.6pt,minimum size=3mm},
  lab/.style={font=\footnotesize\sffamily,text=slate,align=center}
]
\draw[->,rulegray] (-0.4,0) -- (7.2,0);
\draw[thick,rulegray] (-0.2,-0.75) -- (6.8,2.1);
\foreach \x in {-0.2,0.6,1.4,2.2,3.0,3.8,4.6,5.4,6.2,6.8}
  \node[pt] at (\x,{\x*0.35-0.68}) {};
\node[sol] (s0) at (2.2,0.09) {};
\node[sol] (s1) at (3.8,0.65) {};
\node[sol] (s2) at (5.4,1.21) {};
\node[lab] at (2.2,-0.45) {$x_0$};
\node[lab] at (3.8,1.15) {$x_0 + \tfrac{b}{d}$};
\node[lab] at (5.4,1.71) {$x_0 + 2\tfrac{b}{d}$};
\node[lab] at (3.2,-1.25) {Once one solution exists, all solutions are
spaced $\tfrac{b}{\gcd}$ apart --- an arithmetic family,\\
never a scatter. (For $ax + by = c$.)};
\end{tikzpicture}
